% Part: model-theory
% Chapter: interpolation
% Section: definability

\documentclass[../../../include/open-logic-section]{subfiles}

\begin{document}

\olfileid{mod}{int}{def}

\olsection{నిర్వచనీయతా సిద్ధాంతం}

అంతర్వేశ సిద్ధాంతానికి ఒక ముఖ్యమైన అన్వయం బెత్ నిర్వచనీయతా
సిద్ధాంతం. $n$-స్థానాల సంబంధం~$R$ను నిర్వచించడానికి, $R$ రాని,
$n$ స్వేచ్ఛా \tetoken{చరాలు}{!!{variable}s} గల
\tetoken{ఒక సూత్రం}{!!a{formula}}~$!C$ను ఇవ్వవచ్చు. ఇది $!C$ ద్వారా
$R$కు ఇచ్చే \emph{స్పష్ట} నిర్వచనం. అప్పుడు $n$-స్థానాల
\tetoken{విధేయ సంకేతం}{!!{predicate}}~$P$ను కలిగిన
\tetoken{భాష}{!!{language}}లోని సిద్ధాంతం~$\Sigma(P)$, ఒక
లాంఛనీకృత స్పష్ట నిర్వచనాన్ని కలిగి ఉంటే (లేదా కనీసం దాన్ని
అనుగమింపజేస్తే), $P$ను స్పష్టంగా నిర్వచిస్తుంది అని కూడా చెప్పవచ్చు;
అంటే,
\[
\Sigma(P) \Entails \lforall[x_1][\dots
  \lforall[x_n][(\Atom{P}{x_1,\dots, x_n} \liff !C(x_1, \dots,
    x_n))]].
\]
అయితే ఒక సంబంధాన్ని నిర్వచించడానికి---అంటే దాన్ని పూర్తిగా
నిర్ణయించడానికి---స్పష్ట నిర్వచనం ఒక్కటే మార్గం కాదు. ఏ నమూనాలోనైనా
మిగిలిన \tetoken{భాష}{!!{language}}కు ఇచ్చిన అర్థనిర్దేశం వల్లనే
$P$ యొక్క అర్థనిర్దేశం స్థిరపడేలా కూడా ఒక సిద్ధాంతం ఉండవచ్చు. ఒక
సిద్ధాంతం $P$ అర్థనిర్దేశాన్ని ఇలా స్థిరపరిచినప్పుడల్లా---అంటే $P$ను
\emph{అవ్యక్తంగా నిర్వచించినప్పుడల్లా}---దాన్ని స్పష్టంగా కూడా
నిర్వచిస్తుందని నిర్వచనీయతా సిద్ధాంతం చెబుతుంది.

\begin{defn}
\tetoken{విధేయ సంకేతం}{!!{predicate}}~$P$ను కలిగి లేని
\tetoken{ఒక భాష}{!!a{language}}~$\Lang{L}$ అనుకుందాం.
$\Lang{L} \cup \{P\}$లోని \tetoken{వాక్యాల}{!!{sentence}s}
సమితి~$\Sigma(P)$, $P$ను \emph{స్పష్టంగా నిర్వచిస్తుంది} అంటే,
కింది షరతు తీరే $\Lang{L}$కు చెందిన
\tetoken{ఒక సూత్రం}{!!a{formula}}~$!C(x_1, \dots, x_n)$ ఉంటుంది:
\[
\Sigma(P) \Entails \lforall[x_1][\dots
  \lforall[x_n][(\Atom{P}{x_1,\dots, x_n} \liff !C(x_1, \dots,
    x_n))]].
\]
\end{defn}

\begin{defn}
\tetoken{విధేయ సంకేతాలు}{!!{predicate}s}~$P$, $P'$లను కలిగి లేని
\tetoken{ఒక భాష}{!!a{language}}~$\Lang{L}$ అనుకుందాం.
$\Lang{L} \cup \{P\}$లోని \tetoken{వాక్యాల}{!!{sentence}s}
సమితి~$\Sigma(P)$, $P$ను \emph{అవ్యక్తంగా నిర్వచిస్తుంది} అంటే,
\[
\Sigma(P) \cup \Sigma(P') \Entails \lforall[x_1][\dots
  \lforall[x_n][(\Atom{P}{x_1,\dots, x_n} \liff \Atom{P'}{x_1,\dots,
      x_n})]],
\]
ఇక్కడ $\Sigma(P)$ అంతటా $P$ స్థానంలో $P'$ను ఏకరీతిగా
ప్రతిస్థాపించడం వల్ల వచ్చేదే $\Sigma(P')$.
\end{defn}

మరో మాటలో చెప్పాలంటే, ఏ నమూనా $\Struct{M}$కైనా,
$R, R' \subseteq \Domain{M}^n$కైనా, $\Expan{M}{R} \Entails
\Sigma(P)$, $\Expan{M}{R'} \Entails \Sigma(P')$ రెండూ అయితే
$R=R'$. ఇక్కడ $\Expan{M}{R}$ అనేది $\Lang{L}$ను
$\Lang{L} \cup \{P\}$కు విస్తరించిన
\tetoken{నిర్మాణం}{!!{structure}}~$\Struct{M'}$; అందులో
$\Assign{P}{M'} = R$. $\Expan{M}{R'}$ కూడా ఇలాగే ఉంటుంది.

\begin{thm}[బెత్ నిర్వచనీయతా సిద్ధాంతం]
$\Lang{L} \cup\{P\}$-
\tetoken{సూత్రాల}{!!{formula}s} సమితి~$\Sigma(P)$, $P$ను అవ్యక్తంగా
నిర్వచించినప్పుడు, అప్పుడు మాత్రమే $\Sigma(P)$, $P$ను స్పష్టంగా
నిర్వచిస్తుంది.
\sourcecorrection{OLTEMODINT-007}{ఆంగ్ల మూల సిద్ధాంత వాక్యంలో
“if and only” తరువాత కావాల్సిన రెండవ “if” తప్పిపోయింది. ముందరి రెండు
నిర్వచనాలు, తరువాతి రెండు దిశల నిరూపణ కోరే ద్విశరతను తెలుగులో
పూర్తిగా పునరుద్ధరించాం.}
\end{thm}

\begin{proof}
$\Sigma(P)$, $P$ను స్పష్టంగా నిర్వచిస్తే కింది రెండూ తీరుతాయి:
\begin{align*}
  \Sigma(P) & \Entails & \lforall[x_1][\dots \lforall[x_n]
    [(\Atom{P}{x_1,\dots, x_n} \liff !C(x_1,\dots,x_n))]]\\
  \Sigma(P') & \Entails & \lforall[x_1][\dots \lforall[x_n]
    [(\Atom{P'}{x_1,\dots, x_n} \liff !C(x_1,\dots,x_n))]]
\end{align*}
అందువల్ల కావాల్సిన నిష్కర్ష వస్తుంది. తిరుగు దిశ కోసం $\Sigma(P)$,
$P$ను అవ్యక్తంగా నిర్వచిస్తుందని అనుకుందాం. మొదట
\tetoken{స్థిరసంకేతాలు}{!!{constant}s}~$c_1$, \dots,~$c_n$ను
$\Lang{L}$కు చేర్చుతాం. అప్పుడు
\[
\Sigma(P) \cup \Sigma(P') \Entails
\Atom{P}{c_1, \dots, c_n} \to  \Atom{P'}{c_1, \dots, c_n}.
\]
సంహతత్వం వల్ల $\Delta_0 \subseteq \Sigma(P)$,
$\Delta_1 \subseteq \Sigma(P')$ అనే పరిమిత సమితులు ఉండి,
\[
\Delta_0 \cup \Delta_1 \Entails
\Atom{P}{c_1, \dots, c_n} \to \Atom{P'}{c_1, \dots, c_n}.
\]
$!A(P) \in \Delta_0$ లేదా $!A(P') \in \Delta_1$ అయ్యే అన్ని
\tetoken{వాక్యాల}{!!{sentence}s}~$!A(P)$ సంయోగాన్ని $!D(P)$గా;
$!A(P) \in \Delta_0$ లేదా $!A(P') \in \Delta_1$ అయ్యే అన్ని
\tetoken{వాక్యాల}{!!{sentence}s}~$!A(P')$ సంయోగాన్ని $!D(P')$గా
ఉంచండి. అప్పుడు $!D(P) \land !D(P') \Entails
\Atom{P}{c_1, \dots, c_n} \to \Atom{P'}{c_1,\dots,c_n}$.
\sourcecorrection{OLTEMODINT-008}{ఈ అనుగమనపు చివరి పరమాణు
సూత్రాన్ని ఆంగ్ల మూలం $P'c_1\dots c_n$గా ముడి సంకేతాల వరుసగా
రాస్తుంది; అదే అనుగమనంలో ముందే వాడిన, వ్యాకరణపరంగా ఏర్పడిన
$\Atom{P'}{c_1,\dots,c_n}$ రూపాన్ని పునరుద్ధరించాం.}
ప్రతి \tetoken{విధేయ సంకేతం}{!!{predicate}} $\Entails$కు ఒక్కవైపే
వచ్చేలా దీన్ని పునర్వ్యవస్థీకరించవచ్చు:
\[
!D(P) \land \Atom{P}{c_1, \dots, c_n} \Entails
!D(P') \to \Atom{P'}{c_1, \dots, c_n}.
\]
క్రేగ్ అంతర్వేశ సిద్ధాంతం వల్ల $P$గానీ $P'$గానీ లేని
\tetoken{ఒక వాక్యం}{!!a{sentence}}~$!C(c_1,\dots, c_n)$ ఉండి:
\begin{align*}
  !D(P) \land \Atom{P}{c_1, \dots, c_n} & \Entails !C(c_1,\dots, c_n); \\
  !C(c_1,\dots, c_n) & \Entails !D(P') \to \Atom{P'}{c_1, \dots, c_n}.
\end{align*}
ఈ రెండు అనుగమనాల్లో మొదటిదాని నుంచి $!D(P) \Entails
\Atom{P}{c_1,\dots, c_n} \lif !C(c_1,\dots, c_n)$ వస్తుంది.
రెండవదాని విషయంలో, $\Lang{L} \cup \{P\}$-నమూనా
$\Expan{M}{R} \Entails !A(P)$ అయినప్పుడు, అప్పుడు మాత్రమే అనురూప
$\Lang{L} \cup \{P'\}$-నమూనా $\Expan{M}{R} \models !A(P')$ కాబట్టి,
$!C(c_1,\dots, c_n) \Entails !D(P) \lif
\Atom{P}{c_1,\dots, c_n}$ వస్తుంది. దీని నుంచి
\[
!D(P) \Entails !C(c_1,\dots,c_n) \to \Atom{P}{c_1,\dots, c_n}.
\]
ఈ రెండింటినీ కలిపితే $!D(P) \Entails \Atom{P}{c_1,\dots, c_n}
\liff !C(c_1, \dots, c_n)$; ఏకదిశత, సార్వత్రీకరణ ద్వారా
\[
\Sigma(P) \Entails
\lforall[x_1][\dots\lforall[x_n][(\Atom{P}{x_1,\dots, x_n} \liff
    !C(x_1,\dots, x_n))]].
\]
కూడా వస్తుంది.
\end{proof}

\end{document}
